\providecommand{\kBKM}{\kappa_{\mathrm{BKM}}}
\providecommand{\kcat}{\kappa_{\mathrm{cat}}}
\providecommand{\kch}{\kappa_{\mathrm{ch}}}
\providecommand{\kfib}{\kappa_{\mathrm{fiber}}}
\providecommand{\Z}{\mathbb{Z}}
\providecommand{\DeltaFive}{\Delta_5}
\providecommand{\DFive}{D_5}
\providecommand{\fg}{\mathfrak{g}}

\section*{Rank-Three Paramodular Ladder at
$N\in\{2,3,4,6\}$}
\label{sec:gn-rank3-bkm-tower}

The primitive CHL denominator ladder contains the five levels
\[
  N\in\{1,2,3,4,6\}.
\]
At level \(1\), Gritsenko--Nikulin construct the rank-three
Borcherds--Kac--Moody superalgebra \(\fg_{\DeltaFive}\) with denominator
\(\DeltaFive\) (Gritsenko--Nikulin 1997, 1998). At levels \(N\geq2\),
the constructed object in the cited sources is the paramodular
Borcherds product \(\Phi_N\). A Lie-superalgebra theorem at these
levels requires a Lorentzian root lattice, real simple roots, Weyl
vector, parity decomposition, root supermultiplicity function, and
Weyl--Borcherds denominator identity.

\subsection*{Borcherds Constants}

Let \(\phi^{(N)}_{0,1}\) be the \(g_N\)-twined weak Jacobi form of
weight \(0\) and index \(1\) used as the arithmetic-lift input for the
level-\(N\) CHL denominator. Let \(c_N(0)\) denote its discriminant-zero
constant term in the Borcherds-product normalization.
Gritsenko--Nikulin's \(\Delta\)-tower, Gritsenko's arithmetic lift, and
Borcherds' weight theorem give
\[
  (c_1(0),c_2(0),c_3(0),c_4(0),c_6(0))
  =
  (10,8,6,4,2),
\]
and the paramodular weights are
\[
  \operatorname{wt}(\Phi_N)
  =
  \frac{c_N(0)}{2}
  =
  (5,4,3,2,1)
\]
for \(N=(1,2,3,4,6)\).

\begin{proof}
For \(N=1\), the Eichler--Zagier generator satisfies
\[
  \phi_{0,1}(\tau,z)|_{q^0}
  =
  \zeta^{-1}+10+\zeta ,
\]
so \(c_1(0)=10\). Gritsenko's level-one arithmetic lift sends this input
to the weight-\(5\) paramodular form \(\DeltaFive\) (Gritsenko 1999,
Theorem~3.4), and Borcherds' product theorem gives
\[
  \operatorname{wt}(\DeltaFive)=c_1(0)/2=5.
\]

For \(N=2,3,4,6\), the Gritsenko--Nikulin CHL tower supplies
paramodular cusp forms
\[
  \Delta_{4}^{(2)},\qquad
  \Delta_{3}^{(3)},\qquad
  \Delta_{2}^{(4)},\qquad
  \Delta_{1}^{(6)}
\]
of weights \(4,3,2,1\), with direct product checks for the \(N=4\)
and \(N=6\) entries in Cl\'ery--Gritsenko 2013 and
Cl\'ery--Gritsenko--Hulek 2015. Borcherds 1998, Theorem~13.3,
identifies the weight with \(c_N(0)/2\). Hence
\[
  c_2(0)=8,\qquad c_3(0)=6,\qquad c_4(0)=4,\qquad c_6(0)=2 .
\]
The same table is obtained from the twined K3 elliptic genus in the
standard CHL normalization. Govindarajan--Krishna 2010, Table~1, list
the square partition-function weights
\(2\operatorname{wt}(\Phi_N)=(10,8,6,4,2)\).
\end{proof}

The CHL ladder and the Gritsenko--Cl\'ery diagonal-divisor catalogue are
distinct Borcherds-product families. The Gritsenko--Cl\'ery twice-weight
tuple is
\[
  (10,4,6,2,4,1,3,2),
\]
and the only canonical common entry is the level-one form
\(\Delta_5\). At \(N=6\), the CHL denominator belongs to the
Cl\'ery--Gritsenko--Hulek Niemeier-\(6D_4\) construction. No row map
identifies it with the \(A_5^4D_4\) Niemeier entry.

Thus the BKM modular characteristic on this ladder is
\[
  \kBKM(\Phi_N)=c_N(0)/2 .
\]
This invariant records the weight of the denominator product. It is
separate from \(\kappa_{\mathrm{ch}}\), \(\kappa_{\mathrm{cat}}\), and
\(\kappa_{\mathrm{fiber}}\).

\subsection*{\texorpdfstring{$K3\times E$}{K3 x E} Invariant Values}

For \(X=K3\times E\), the compact categorical and chiral supertraces are
direct computations:
\[
  \kcat(X)=\chi(\mathcal{O}_{K3})\chi(\mathcal{O}_E)=2\cdot0=0,
  \qquad
  \kch(A_X)=\sum_{q=0}^{3}(-1)^q h^{0,q}(X)=1-1+1-1=0.
\]
The Heisenberg--Mukai specialization, the Borcherds denominator, and the
K3 Mukai lattice give the remaining three values:
\[
  \kch^{\mathrm{Heis}}(X)=3,\qquad
  \kBKM(\fg_{\DeltaFive})=5,\qquad
  \kfib(K3)=24.
\]
The value \(\kBKM(\fg_{\DeltaFive})=5\) is the \(N=1\) case of
\(\kBKM(\Phi_N)=c_N(0)/2\). The \(d=3\) native output of
\(\Phi\) is \(E_1\)-chiral. The denominator construction belongs to the
Hall--Drinfeld/BKM branch. The Mukai self-mirror K3 current branch uses
separate stable-envelope data.

\subsection*{The Level-One Lie Superalgebra}

The complete rank-three BKM construction supplied by Gritsenko--Nikulin
is the level-one algebra \(\fg_{\DeltaFive}\). Its reflection lattice
has basis \(f_2,f_3,f_{-2}\) with
\[
  (f_2,f_{-2})=-1,\qquad (f_3,f_3)=2,
\]
and real simple roots
\[
  \delta_1=2f_2-f_3,\qquad
  \delta_2=2f_{-2}-f_3,\qquad
  \delta_3=f_3 .
\]
Their Gram matrix is
\[
  \begin{pmatrix}
    2&-2&-2\\
    -2&2&-2\\
    -2&-2&2
  \end{pmatrix},
  \qquad \det=-32,
\]
of signature \((2,1)\). The Weyl vector is
\[
  \rho=\frac12(\delta_1+\delta_2+\delta_3)
  =f_2-\frac12 f_3+f_{-2},
  \qquad
  (\rho,\delta_i)=-1 .
\]
Set
\[
  \DFive(Z)=64^{-1}\DeltaFive(Z).
\]
The monic denominator section is
\[
  \operatorname{den}(\fg_{\DeltaFive})(Z)=\DFive(2Z)
  =64^{-1}\DeltaFive(2Z),
\]
and its product exponents are the signed root supermultiplicities read
from the Fourier coefficients of \(\phi_{0,1}\) by the discriminant
\(4nm-\ell^2\). The real simple roots have multiplicity \(1\) as part
of the Cartan datum.

\subsection*{Higher-Level Root Data}

For \(N\in\{2,3,4,6\}\), a rank-three BKM superalgebra theorem requires
the datum
\[
  (L_N,\Delta^{\mathrm{re}}_N,\rho_N,\epsilon_N,m_N,\Phi_N).
\]
Here \(L_N\) is a Lorentzian lattice of signature \((2,1)\),
\(\Delta^{\mathrm{re}}_N\) is a set of real simple roots,
\(\rho_N\) is a Weyl vector, \(\epsilon_N\) is the parity rule,
\(m_N\) is the full root-multiplicity function, and \(\Phi_N\) is the
Weyl--Borcherds denominator attached to these choices. The known
paramodular form \(\Phi_N\) supplies the final automorphic factor and
its weight. The Lie bracket and the parity of imaginary simple roots
require the remaining root-space data.

The common construction across the five levels is the level-dependent
Gritsenko arithmetic lift
\[
  \mathrm{Lift}_N:\phi^{(N)}_{0,1}\longmapsto \Phi_N .
\]
The Weyl groups, root lattices, and characters are outputs of the
level-\(N\) construction. The cited sources give no single common Weyl
group for all five levels.

\subsection*{The \(N=6\) Lattice and Quotient Data}

The \(N=6\) CHL/umbral anchor is the Niemeier lattice with root system
\[
  6D_4 .
\]
The lattice \(4A_5+D_4\) is a different Niemeier lattice of Coxeter
number \(6\), with a different umbral character from the \(6A\) CHL
denominator.

The cyclic quotient singularities of \(K3/\langle g_6\rangle\) form a
third datum. Nikulin's and Xiao's fixed-point counts give
\[
  \#\operatorname{Fix}(g_6)=2,\qquad
  \#\operatorname{Fix}(g_6^2)=6,\qquad
  \#\operatorname{Fix}(g_6^3)=8.
\]
The quotient has singularity type
\[
  2A_5+2A_2+2A_1,
\]
and the Euler check is
\[
  e(K3/\Z_6)=\frac{24+2+6+8+6+2}{6}=8,\qquad
  8+(2\cdot5+2\cdot2+2\cdot1)=24 .
\]
Niemeier anchors are global positive-definite even unimodular lattices.
Du Val singularities are local analytic quotient germs. Their ADE
labels encode different data.

\subsection*{First Coefficient Check at \(N=2\)}

At \(N=2\), the CHL constant is \(c_2(0)=8\), so
\[
  \operatorname{wt}(\Phi_2)=4 .
\]
This verifies the first denominator constant. A real-root multiplicity
statement requires the Cartan lattice, Weyl vector, parity, and
multiplicity rule.

\subsection*{Conclusion}

The theorem-level statement is the paramodular one:
\[
  \Phi_N \text{ exists and } \kBKM(\Phi_N)=c_N(0)/2
  \quad (N\in\{1,2,3,4,6\}).
\]
The level-one case further carries the constructed BKM superalgebra
\(\fg_{\DeltaFive}\). Higher-level Lie-superalgebra statements require
the root datum \((L_N,\Delta^{\mathrm{re}}_N,\rho_N,\epsilon_N,m_N)\).
The \(N=6\) geometry uses the \(6D_4\) Niemeier anchor and the cyclic
quotient type \(2A_5+2A_2+2A_1\); \(4A_5+D_4\) belongs to a different
Niemeier entry.
